% Comment ce livre a été fabriqué, et le tableau des vingt contributions.
\chapter{Comment ce livre a été fabriqué}
\markboth{Comment ce livre a été fabriqué}{Comment ce livre a été fabriqué}

\section{Qui a écrit quoi}

Ce livre a été écrit par des agents d'IA (Claude, d'Anthropic) sous la direction de l'auteur. Un agent a conçu le
cadre : la mise en page, le prologue, le chapitre sur la méthode, le chapitre de synthèse, la lettre et ces pages de
fin. Quatre autres agents ont écrit en parallèle les vingt chapitres de contribution, à partir d'un même guide de
style et d'une même fiche par contribution, et avec une règle stricte : n'utiliser aucun fait ni aucun nombre qui ne
figure pas dans les dossiers des travaux eux-mêmes. L'auteur a choisi le sujet, le ton, la dédicace, et il répond du
contenu. Au moment où ces lignes sont écrites, aucun scientifique extérieur n'a relu ce livre.

Les résultats racontés ont été produits de la même façon, par des agents d'IA sous la direction de l'auteur, et
audités par d'autres agents d'IA, selon la méthode décrite au chapitre~\ref{ch:elenchus}. \textbf{Aucun mathématicien
humain ne les a relus.}

\section{Comment les vingt ont été choisis}

\paragraph{L'inventaire} Tous les dépôts publiés sur Zenodo pendant l'année ont été recensés : 144 dépôts, qui
correspondent à 112 travaux distincts, une fois fusionnés les doublons (le même texte déposé jusqu'à treize fois) et
les versions successives d'un même travail. Le compte GitHub de l'auteur compte 420 dépôts, dont 317 copies de
projets d'autres personnes ; 82 sont des projets propres à l'auteur, créés ou modifiés pendant l'année (de fin
octobre 2025 à fin octobre 2026), et 57 d'entre eux touchent à la science. Un dépôt GitHub n'est retenu comme candidat
que si aucun travail publié sur Zenodo ne le décrit déjà. Au total, 113 candidats.

\paragraph{La grille, fixée avant de noter} La grille de sélection a été écrite et gelée le 26 octobre 2026,
\emph{avant} qu'aucun candidat ne soit noté, et elle n'a pas été modifiée ensuite. Elle note chaque candidat de 0 à 3
sur sept critères, la même grille pour tous, quel que soit le domaine : la rigueur atteinte (le niveau Elenchus le plus
élevé), la vérification indépendante, la reproductibilité, la synthèse entre mathématiques, physique ou biologie,
calcul et preuve formelle, l'honnêteté (résultats négatifs, corrections publiées), la nouveauté, et la valeur pour
des non-spécialistes. Chaque note est justifiée en une ligne, avec le fichier sur lequel elle repose.

\paragraph{Les règles fixes} Quatre règles, appliquées dans cet ordre, complètent la grille. Les travaux dont
l'affirmation centrale a été retirée ou réfutée ne sont pas notés comme des succès : ils sont rattachés au travail
qui publie le retrait, et notés comme des leçons. Deux places au moins sont réservées à chacun des cinq domaines
(mathématiques, physique, génétique et biologie, calcul numérique et logiciel, méthode et vérification de l'IA). Au
moins trois des vingt doivent être des leçons. Aucun domaine ne peut avoir plus de sept places. Les domaines sont
ceux de l'inventaire : la cosmologie et l'enseignement y sont comptés à part de la physique, ce qui explique neuf
chapitres de physique au sens large (sept de physique et deux de cosmologie).

\paragraph{Ce que les règles ont changé} Par rapport aux vingt meilleures notes brutes, les règles ont fait entrer
quatre contributions (l'instrument rusty-SUNDIALS, la porte GWAYA, l'établi AutoevolveAI et la leçon du stellarator)
pour respecter les minimums des domaines « calcul numérique » et « méthode », et en ont fait sortir quatre : deux pour
le plafond de sept en physique, deux déplacées par les places réservées. Dix-sept des vingt contributions datent
du 8 au 26 octobre 2026 : la grille récompense le rejeu, les audits indépendants avec corrections publiées et les résultats
négatifs, que le projet n'a produits de façon systématique qu'à partir d'octobre. Les mois précédents entrent dans le
livre surtout par les leçons. Le dossier complet de la sélection, avec toutes les notes et leurs justifications, est
public.\reference{Dossier \nolinkurl{results/discovery_journal_2026/} : \texttt{RUBRIC.md},
\nolinkurl{SELECTION_final.md}, \nolinkurl{TOP20_final.json}, \nolinkurl{scores_final.json}, \texttt{RESCORE.md}.}

\paragraph{Une note revue en cours de route} La première notation s'appuyait sur le registre de rejeu tel qu'il
était à mi-parcours (71 modules rejoués). Quand le rejeu s'est terminé (93 modules, tous passés), les notes qui
citaient les modules nouvellement rejoués ont été recalculées, avec la même grille, sans la modifier. Une seule
contribution de la liste a changé : le certificat sur les forces écrantées est sorti, à égalité de note, départagé
par la date puis écarté par le plafond de sept en physique, et la réfutation numérique de la conjecture QU-1
(chapitre~\ref{ch:c12}) est entrée. L'ordre de quelques autres a bougé. Ce livre suit la liste finale ; l'ancienne
et la nouvelle sont toutes deux conservées dans le dossier.

\section{Les vingt contributions}

Le tableau donne, pour chaque contribution, son rang, le titre de son chapitre, son domaine, sa note de rigueur dans
la grille et son DOI. La note de rigueur mesure l'affirmation la mieux vérifiée que le travail énonce lui-même : 3
pour une preuve vérifiée par le juge et rejouée, ou pour un certificat de calcul recontrôlé par un programme
différent ; 2 pour une preuve vérifiée par le juge avec audit des axiomes, sans rejeu ; 1 pour un argument, une
citation ou un calcul exact sans le juge ; 0 pour des nombres à virgule seulement. Ce n'est pas le niveau de
l'ensemble du travail, et ce n'est pas tout à fait un niveau Elenchus : les pastilles de niveau, morceau par morceau,
se trouvent dans chaque chapitre, et elles font foi. (La note de la contribution n°\,20 suit l'inventaire, qui ne
cite pas la liste des axiomes ; le fichier de contrôle du dépôt la donne, d'où la pastille As du chapitre.)

\begingroup
\small
\renewcommand{\arraystretch}{1.22}
% generated by code/make_backmatter.py from TOP20_final.json and chapters/ORDER.json -- do not edit
\begin{tabularx}{\linewidth}{@{}r>{\raggedright\arraybackslash}X>{\raggedright\arraybackslash}p{2.3cm}c@{}}
\toprule
\textsf{n°} & \textsf{contribution (titre du chapitre), DOI ou dépôt} & \textsf{domaine} & \textsf{rigueur}\\
\midrule
1 & Des distances cosmiques avec une marge d'erreur vérifiée par une machine \newline {\scriptsize\color{encre!60}2026-10-08, p.~\pageref{ch:c01}} \newline{\footnotesize \doi{10.5281/zenodo.23247730}} & cosmologie & {\sffamily\small 3\,/\,3}\\
2 & Une tension célèbre de la cosmologie, écrite comme une inégalité vérifiée \newline {\scriptsize\color{encre!60}2026-10-10, p.~\pageref{ch:c02}} \newline{\footnotesize \doi{10.5281/zenodo.23284020}} & cosmologie & {\sffamily\small 3\,/\,3}\\
3 & Le triangle, rangement le moins coûteux : des conséquences vérifiées et une simulation \newline {\scriptsize\color{encre!60}2026-10-11, p.~\pageref{ch:c03}} \newline{\footnotesize \doi{10.5281/zenodo.23293429}} & physique & {\sffamily\small 3\,/\,3}\\
4 & La répulsion écrantée entre dans la famille du théorème, et une erreur est réparée \newline {\scriptsize\color{encre!60}2026-10-25, p.~\pageref{ch:c04}} \newline{\footnotesize \doi{10.5281/zenodo.23646105}} & physique & {\sffamily\small 3\,/\,3}\\
5 & Refaire le calcul d'un autre, sans le déclarer juste \newline {\scriptsize\color{encre!60}2026-10-26, p.~\pageref{ch:c13}} \newline{\footnotesize \doi{10.5281/zenodo.23653827}} & mathématiques & {\sffamily\small 3\,/\,3}\\
6 & Une règle prouvée droite montre que notre propre esquisse était fausse \newline {\scriptsize\color{encre!60}2026-10-26, p.~\pageref{ch:c05}} \newline{\footnotesize \doi{10.5281/zenodo.23653846}} & physique & {\sffamily\small 3\,/\,3}\\
7 & Quelle chance a une mutation d'envahir une population~? \newline {\scriptsize\color{encre!60}2026-10-26, p.~\pageref{ch:c06}} \newline{\footnotesize \doi{10.5281/zenodo.23653850}} & génétique, biologie & {\sffamily\small 3\,/\,3}\\
8 & Quand le plan de table ne change rien à la sélection \newline {\scriptsize\color{encre!60}2026-10-26, p.~\pageref{ch:c07}} \newline{\footnotesize \doi{10.5281/zenodo.23660071}} & génétique, biologie & {\sffamily\small 3\,/\,3}\\
9 & Un manuel pour suivre un théorème jusqu'au laboratoire \newline {\scriptsize\color{encre!60}2026-10-11, p.~\pageref{ch:c08}} \newline{\footnotesize \doi{10.5281/zenodo.23295142}} & enseignement & {\sffamily\small 3\,/\,3}\\
10 & Des paires électron--trou qui se rangent en triangles, sur le papier \newline {\scriptsize\color{encre!60}2026-10-25, p.~\pageref{ch:c09}} \newline{\footnotesize \doi{10.5281/zenodo.23646103}} & physique & {\sffamily\small 3\,/\,3}\\
\bottomrule
\end{tabularx}

\bigskip

\begin{tabularx}{\linewidth}{@{}r>{\raggedright\arraybackslash}X>{\raggedright\arraybackslash}p{2.3cm}c@{}}
\toprule
\textsf{n°} & \textsf{contribution (titre du chapitre), DOI ou dépôt} & \textsf{domaine} & \textsf{rigueur}\\
\midrule
11 & La force dans un semi-conducteur atomiquement mince, en tas de sable \newline {\scriptsize\color{encre!60}2026-10-25, p.~\pageref{ch:c10}} \newline{\footnotesize \doi{10.5281/zenodo.23646107}} & physique & {\sffamily\small 3\,/\,3}\\
12 & Tester avant de prouver : une idée naturelle qui ne tient pas \emph{(leçon)} \newline {\scriptsize\color{encre!60}2026-10-26, p.~\pageref{ch:c12}} \newline{\footnotesize \doi{10.5281/zenodo.23653833}} & physique & {\sffamily\small 3\,/\,3}\\
13 & Parfait au repos, pas quand il frissonne \emph{(leçon)} \newline {\scriptsize\color{encre!60}2026-10-26, p.~\pageref{ch:c11}} \newline{\footnotesize \doi{10.5281/zenodo.23655392}} & physique & {\sffamily\small 3\,/\,3}\\
14 & Quatre étages sur douze, et le prix du reste \newline {\scriptsize\color{encre!60}2026-10-26, p.~\pageref{ch:c14}} \newline{\footnotesize \doi{10.5281/zenodo.23656804}} & mathématiques & {\sffamily\small 3\,/\,3}\\
15 & Remonter la machine des autres dans son propre atelier \newline {\scriptsize\color{encre!60}2026-10-26, p.~\pageref{ch:c15}} \newline{\footnotesize \doi{10.5281/zenodo.23661381}} & mathématiques & {\sffamily\small 3\,/\,3}\\
16 & Une garantie générale, trois certificats précis \newline {\scriptsize\color{encre!60}2026-10-10, p.~\pageref{ch:c16}} \newline{\footnotesize \doi{10.5281/zenodo.23285631}} & mathématiques & {\sffamily\small 3\,/\,3}\\
17 & Le moteur des calculs, et son rapport d'inspection \newline {\scriptsize\color{encre!60}2026-05-11, p.~\pageref{ch:c17}} \newline{\footnotesize dépôt \texttt{rusty-SUNDIALS}} & calcul numérique & {\sffamily\small 0\,/\,3}\\
18 & Une porte qui reste fermée quand on ne sait pas \newline {\scriptsize\color{encre!60}2026-10-03, p.~\pageref{ch:c18}} \newline{\footnotesize \doi{10.5281/zenodo.23119953}} & méthode, vérification & {\sffamily\small 2\,/\,3}\\
19 & Le banc d'essai qui s'est inspecté lui-même \newline {\scriptsize\color{encre!60}2026-09-19, p.~\pageref{ch:c19}} \newline{\footnotesize dépôt \texttt{AutoevolveAI}} & méthode, vérification & {\sffamily\small 1\,/\,3}\\
20 & On accusait le thermostat de surchauffer \emph{(leçon)} \newline {\scriptsize\color{encre!60}2026-10-21, p.~\pageref{ch:c20}} \newline{\footnotesize \doi{10.5281/zenodo.23532458}} & calcul numérique & {\sffamily\small 1\,/\,3}\\
\bottomrule
\end{tabularx}

\endgroup
